\documentclass[12pt]{article}

\usepackage[a4paper, margin=1in]{geometry}
\usepackage{amsmath,amssymb,amsthm}
\usepackage{mathtools}
\usepackage{mathrsfs}
\usepackage{array}
\usepackage{booktabs}
\usepackage{longtable}
\usepackage{tabularx}
\usepackage{enumitem}
\usepackage{float}
\usepackage{xcolor}
\usepackage[hidelinks]{hyperref}
\usepackage{fontspec}
\usepackage{xeCJK}

\defaultfontfeatures{Ligatures=TeX}
\IfFontExistsTF{Times New Roman}
  {\setmainfont{Times New Roman}}
  {\setmainfont{TeX Gyre Termes}}
\IfFontExistsTF{Menlo}
  {\setmonofont{Menlo}}
  {\setmonofont{Latin Modern Mono}}
\IfFontExistsTF{PingFang SC}
  {\setCJKmonofont{PingFang SC}}
  {\IfFontExistsTF{Songti SC}
    {\setCJKmonofont{Songti SC}}
    {}}
\IfFontExistsTF{PingFang SC}
  {\setCJKmainfont{PingFang SC}}
  {\IfFontExistsTF{Songti SC}
    {\setCJKmainfont{Songti SC}}
    {\setCJKmainfont{STSong}}}

\newtheorem{theorem}{Theorem}[section]
\newtheorem{proposition}[theorem]{Proposition}
\newtheorem{lemma}[theorem]{Lemma}
\newtheorem{corollary}[theorem]{Corollary}
\theoremstyle{definition}
\newtheorem{definition}[theorem]{Definition}
\newtheorem{example}[theorem]{Example}
\theoremstyle{remark}
\newtheorem{remark}[theorem]{Remark}

\newcolumntype{P}[1]{>{\raggedright\arraybackslash}p{#1}}
\newcolumntype{Y}{X}

\newcommand{\tightlist}{%
  \setlength{\itemsep}{0pt}\setlength{\parskip}{0pt}}

\newcounter{none}

\setlength{\LTleft}{0pt}
\setlength{\LTright}{0pt}

\title{%
  \textbf{The Inseparability of Geometry and Modular Flow:}\\
  \textbf{A Structural Inseparability Theorem and Its Implications for the Hard Problem}%
}

\author{
  Sidong Liu, PhD \\
  \small iBioStratix Ltd \\
  \small \texttt{sidongliu@hotmail.com}
}

\date{April 2026}

\begin{document}
\emergencystretch=2em
\raggedbottom
\hbadness=5000
\vbadness=5000

\maketitle

\begin{abstract}
%
The hard problem of consciousness presupposes that physical structure
and experiential structure can be specified independently. We show that
this presupposition is not available in a precise algebraic setting.

Working with Type III\(_1\) von Neumann algebras, we identify emergent
geometry as the relevant mathematical proxy for physical structure and
Tomita-Takesaki modular flow as the relevant proxy for first-person
temporal structure. The resulting question is whether these two sectors
can be independently deformed within an admissible algebraic class. We
show first that, working within the admissible class selected by
\((U3)\), from the proof-relevant axiom package \[
(U1)+(C)+(Amb)+(CN_\exists),
\] one obtains a standard half-sided modular inclusion and hence, via
the Guido-Longo-Wiesbrock construction, a canonical conformal skeleton.
This yields a conformal-level Weak Inseparability theorem. We then show
that the correct non-gauge local modulus is the double-interval
cross-ratio family, which determines a canonical modular scale datum \[
\{K_b\}_{b\in(0,1)}.
\] Under the additional conditions \((LImp^1)\), \((CR\text{-}nd)\), and
\((SFa)\), this datum induces a Type III-compatible QFIM scale shadow \[
h_b=\mathbf{Q}_{\omega,b}(\partial_b K_b,\partial_b K_b)\,db^2,
\] At the abstract level, this yields a full geometric datum consisting
of the conformal skeleton together with the scale family. In the worked
examples, where the cross-ratio parameter naturally parametrizes spatial
position, this datum determines the Lorentzian metric. In this setting,
the full geometric datum and the extended modular data are
bidirectionally inseparable.

We verify the additional Phase 2 conditions in two structurally distinct
models, the Rindler/free-field chiral half-line model and the Virasoro
net. Both yield the universal profile \[
h_b=C(1-b)^{-2}db^2,
\] with model dependence entering only through the overall constant. The
philosophical consequence is not that the hard problem has been solved
in the ordinary production sense, but that its standard formulation
rests on a presupposition---independent describability---that is not
available in the relevant algebraic setting. The remaining explanatory gap is
thereby relocated from a gap between independently specified domains to
a structural limit on self-description by observers embedded in
\((\mathcal{A},\omega)\).
 \end{abstract}

\clearpage
\tableofcontents
\clearpage

%
\section{The Hidden Presupposition of the Hard
Problem}\label{the-hidden-presupposition-of-the-hard-problem}

The hard problem of consciousness is usually presented as a question
about production: how do physical processes give rise to subjective
experience? In that form, the problem appears deeper than the familiar
``easy problems'' because those concern discrimination, report,
integration, control, and other functional capacities, whereas the hard
problem concerns the existence of a first-person point of view and the
felt character of experience itself. This contrast, introduced most
influentially by Chalmers~\cite{Chalmers1996}, has shaped almost every major debate in
contemporary philosophy of mind.

Yet the canonical formulation contains a strong presupposition that is
rarely isolated and almost never challenged at the level of formal
structure. The presupposition is that physical structure and
experiential structure can be specified independently. One first gives a
complete account of the physical facts, in third-person terms, and then
asks why those facts are accompanied by experience, or by this
experience rather than another. Without this prior separability, the
production question cannot even be posed in its usual form. The hard
problem therefore depends not only on a gap between two domains, but on
the prior availability of those domains as independently describable.

This presupposition is shared by positions that otherwise disagree
sharply. Functionalism and mainstream physicalism accept it and then
attempt to dissolve the residual gap by appeal to organization, role, or
supervenience. Panpsychism accepts it as well, but seeks to avoid brute
emergence by distributing proto-experiential or experiential properties
more fundamentally. Illusionism also inherits the same starting point:
it grants the independent specification of physical structure and then
argues that what seemed left over on the experiential side is a
cognitive misdescription. Integrated Information Theory takes a
different route, but still treats consciousness as something to be
attached to or read off from an already specified structure by means of
an additional quantity. In all of these cases, the debate concerns how
the bridge should be built. The deeper assumption that there are two
independently specifiable shores is left intact.

The present paper takes aim at that assumption. It does not begin by
proposing a new psychophysical law, a new emergent mechanism, or a new
primitive consciousness quantity. It begins by asking whether the
independent-describability presupposition survives translation into a
precise mathematical framework. If ``physical structure'' is represented
not by an unconstrained intuitive picture of matter but by an emergent
geometric structure, and if ``first-person time'' is represented by
modular flow in the sense of Tomita-Takesaki theory, then the hard
problem can be reformulated as a structural question: can these two
sectors be independently varied inside a single admissible algebraic
setting?

The central claim of this paper is that, under a precise Type III\(_1\)
algebraic framework, they cannot. At the conformal level, emergent
causal geometry and conformal modular data are already mutually rigid.
With three additional Phase 2 inputs, the same becomes true at the full
Lorentzian level: complete geometry and extended modular data cannot be
independently deformed. This does not amount to the claim that the hard
problem has simply been ``solved.'' It amounts to a more exact and, in
one sense, more radical thesis: the standard production form of the hard
problem depends on a presupposition that is not available in the
relevant mathematical setting.

This is why the main theorem of the paper is an inseparability theorem
rather than a reduction theorem. The point is not that one side reduces
to the other, nor that one side causally generates the other, but that
both emerge as two readings of one underlying algebraic organization. In
philosophical terms, the paper argues that the usual mind-body contrast
is not fundamental enough to underwrite the standard hard-problem
question. In mathematical terms, it shows that the admissible
deformations of geometry and modular structure do not factor into two
independent sectors.

The argument unfolds in three stages. First, we translate the relevant
philosophical notions into precise mathematical objects: emergent
geometry for the object-side physical structure, and modular flow for
first-person temporal structure. Second, we prove the Phase 1 result:
working within the admissible class selected by \((U3)\), from the
proof-relevant axiom package \[
(U1)+(C)+(Amb)+(CN_\exists)
\] one canonically obtains a conformal skeleton, together with a
conformal-level inseparability theorem. Third, we explain how the full
metric scale is added by the cross-ratio modular datum and its Type III
QFIM shadow, yielding a conditional full inseparability theorem under
the additional Phase 2 package \[
(\mathrm{LImp}^1)+(\mathrm{CR\text{-}nd})+(\mathrm{SFa}).
\] Two worked examples, one based on the Rindler/free-field chiral
half-line model and one on the Virasoro net, show that the Phase 2
conditions are not merely formal placeholders but are realized in
structurally distinct settings.

The resulting philosophical payoff is twofold. First, the hard problem
is reframed: the primary error lies not in failing to build a bridge,
but in assuming that the bridge must connect independently specifiable
domains. Second, the residual explanatory gap is relocated rather than
denied. Once independent describability fails, the remaining question is
no longer how matter produces experience, but why a finite observer
embedded in a single algebraic structure cannot completely internalize
the relation between structural invariants and lived articulation. That
second issue belongs to the limits of self-description, not to a
production gap between two ontologically detached realms.

The rest of the paper is organized as follows. Section 2 identifies the
mathematical counterparts of physical structure and first-person time,
and formulates the independent-describability issue as a question of
independent deformability. Section 3 states the algebraic framework and
axioms. Section 4 proves the Phase 1 result: from the HAFF-native
package to a canonical conformal skeleton and Weak Inseparability.
Section 5 develops the Phase 2 completion via the cross-ratio
modular scale datum and the full inseparability theorems.
Section 6 identifies abstract sufficient conditions for the chart axiom
$(LImp^1)$: for Diff-covariant nets it is fully discharged, and for
non-Diff-covariant nets the remaining burden is isolated to
transportability and regularity assumptions.
Section 7 verifies the remaining Phase 2 conditions in two worked examples.
Section 8 then returns to
the hard problem directly and argues that the theorem shows its
standard formulation rests on a presupposition not available in the
stated algebraic setting, while relocating the remaining explanatory gap to
the self-referential horizon of finite observers. Section 9 discusses
the broader mathematical and physical significance, and Section 10 states
the explicit non-claims and open problems.

\subsection{Logical Roadmap}\label{subsec:roadmap}

The argument has two phases, each with a clear input-output structure.

\emph{Phase 1} (Sections 3--4). Input: the axiom package
\((U1)+(C)+(Amb)+(CN_\exists)\), working within the admissible class
selected by \((U3)\). Output: a canonical conformal skeleton via the
Guido-Longo-Wiesbrock reconstruction, and a conformal-level Weak
Inseparability theorem (Theorem~4.8). The key new step is that
half-sided modular inclusion is \emph{derived} from the
modular-decoherence scaling axiom \((C)\), not assumed.

\emph{Phase 2} (Section 5). Input: the additional conditions
\((LImp^1)+(CR\text{-}nd)+(SFa)\). Output: a full geometric datum
(conformal skeleton + scale shadow) and bidirectional Full
Inseparability (Theorem~5.7). The key new object is the cross-ratio
modular scale datum \(\{K_b\}\) and its Araki-Hessian QFIM shadow
\(h_b\).

Section 6 provides abstract sufficient conditions for \((LImp^1)\).
Section 7 verifies all Phase 2 conditions in two worked examples.
Sections 8--10 draw philosophical consequences, discuss broader
significance, and state open problems.

\section{Mathematical Identification: Geometry as Physics, Modular Flow
as First-Person
Time}\label{mathematical-identification-geometry-as-physics-modular-flow-as-first-person-time}

The philosophical argument of this paper rests on a mathematical
theorem. For the theorem to bear on the hard problem, the translation
between philosophical concepts and mathematical objects must be made
explicit and defended. This section fixes the three identifications that
underwrite the rest of the paper: emergent geometry as the mathematical
content of physical structure, Tomita-Takesaki modular flow as the
mathematical content of first-person time, and independent deformability
as the mathematical translation of independent describability.

\subsection{Emergent Geometry as Physical
Structure}\label{emergent-geometry-as-physical-structure}

In the standard formulation of the hard problem, ``physical structure''
is left deliberately broad. It may refer to microphysical states,
functional organization, or spacetime geometry, depending on the author.
For the present purposes, the relevant notion of physical structure is
not an arbitrary state description but the emergent spacetime geometry
associated with a quantum-algebraic datum.

The reason for this choice is not merely technical convenience. In the
algebraic approach to quantum field theory, spacetime is not a fixed
background container. It is recovered from the algebraic relations among
observables. The programme initiated by Wiesbrock and developed by
Guido, Longo, and others has shown that, under appropriate conditions,
the inclusion structure of von Neumann algebras determines a conformal
geometry~\cite{Wiesbrock1993,GLW1998}, and in favorable cases a full Lorentzian metric. In this
setting, ``physical structure'' is not something given prior to the
algebra. It is something read off from the algebra.

This is the correct level of description for the hard problem, because
the hard problem is not about the relation between consciousness and
some particular physical substrate such as neurons or silicon. It is
about the relation between the subjective and the objective as
structural categories. If the objective-physical side is itself an
emergent reading of algebraic data, then the question of whether it can
be independently specified becomes a question about the algebraic
structure, not about a particular material realization.

We therefore identify:

\[
\text{physical structure} \longleftrightarrow \mathrm{Geom}(\mathcal{N},\omega),
\]

where \(\mathrm{Geom}\) denotes the emergent geometry extracted from a
family \(\mathcal{N}\) of von Neumann algebras in a state \(\omega\). At
the Phase 1 level, this is a conformal skeleton. At the full level, it
includes the metric scale.

\subsection{Modular Flow as First-Person
Time}\label{modular-flow-as-first-person-time}

The second identification is less familiar in the philosophy of mind but
is equally well motivated from the algebraic side.

Given a von Neumann algebra \(\mathcal{A}\) and a faithful normal state
\(\omega\), the Tomita-Takesaki theorem~\cite{Takesaki2003} produces a canonical
one-parameter automorphism group \(\sigma_t^\omega\) of \(\mathcal{A}\),
the modular automorphism group. This group has three properties that
make it the natural mathematical counterpart of first-person temporal
structure.

First, uniqueness. The modular flow is the unique automorphism group
satisfying the KMS condition at inverse temperature \(\beta = 1\) with
respect to \(\omega\). No external clock or background time parameter is
needed to define it. It is intrinsic to the pair
\((\mathcal{A},\omega)\).

Second, algebraic relativity. Different choices of algebra
\(\mathcal{A}\) within the same ambient structure yield different
modular flows. This is not a defect. It is the mathematical expression
of the fact that different observational perspectives give rise to
different temporal experiences. The Unruh effect is the most familiar
physical instance~\cite{Unruh1976}: an accelerated observer's modular flow differs from
that of an inertial observer, and the difference is experienced as a
thermal bath.

Third, structural inseparability at the operator level. The
Tomita-Takesaki polar decomposition \[
S_\Omega = J_\Omega \Delta_\Omega^{1/2}
\] is unique. The modular operator \(\Delta\) and the modular
conjugation \(J\) are not independently choosable components. One cannot
have \(\Delta\) without \(J\), nor \(J\) without \(\Delta\). Since
\(\Delta^{it}\) generates the modular flow and \(J\) implements the
exchange between the algebra and its commutant, this already hints at
the inseparability that the main theorem will establish at the geometric
level: the temporal structure (\(\Delta\)) and the spatial/duality
structure (\(J\)) are two factors of a single operator.

We therefore identify:

\[
\text{first-person time} \longleftrightarrow \sigma_t^\omega = \Delta_\omega^{it}(\cdot)\Delta_\omega^{-it}.
\]

More precisely, the full modular sector includes not only the
one-parameter flow but also the modular Hamiltonian
\(K_\omega = -\ln\Delta_\omega\), the modular conjugation \(J_\omega\),
and, at the extended level, the local modular scale data \(\{K_b\}\)
introduced in Phase 2.

\subsection{From Independent Describability to Independent
Deformability}\label{from-independent-describability-to-independent-deformability}

With the two identifications in place, the philosophical question can be
stated mathematically.

The hard problem, in its standard form, presupposes that one can give a
complete specification of the physical facts and then ask, as a further
question, why those facts are accompanied by experience. Translated into
the present framework, this becomes: can one specify
\(\mathrm{Geom}(\mathcal{N},\omega)\) without thereby fixing
\(\mathrm{Mod}(\mathcal{N},\omega)\), or conversely?

The natural mathematical formalization of ``independent specification''
is independent deformability. Two structures are independently
describable, in the relevant sense, if and only if there exist
admissible deformations of the underlying algebraic data that change one
while leaving the other invariant. If no such deformation exists, then
specifying one structure automatically constrains the other, and the two
cannot be treated as independently given.

This is a precise, falsifiable mathematical question. It does not depend
on intuitions about zombies, inverted qualia, or conceivability
arguments. It depends on the rigidity properties of the map

\[
(\mathcal{N},\omega) \longmapsto \big(\mathrm{Geom}(\mathcal{N},\omega),\;\mathrm{Mod}(\mathcal{N},\omega)\big)
\]

within the admissible class defined by the axiom package.

The hard problem, so translated, becomes:

\begin{quote}
Does there exist an admissible deformation of \((\mathcal{N},\omega)\)
that changes \(\mathrm{Geom}\) while preserving \(\mathrm{Mod}\), or
vice versa?
\end{quote}

The inseparability theorems of this paper answer: no. At the conformal
level (Phase 1), the answer is unconditional within the axiom class. At
the full scale-sensitive level (Phase 2), the answer is conditional on
three
regularity inputs that are verified in two worked examples.

\subsection{What This Translation Does and Does Not
Claim}\label{what-this-translation-does-and-does-not-claim}

Two caveats are essential.

First, the identification of modular flow with first-person time is a
structural correspondence, not an identity claim. We do not assert that
modular flow ``is'' consciousness in any reductive sense. We assert that
modular flow is the correct mathematical object that plays the
structural role assigned to first-person temporal experience in the
hard-problem debate. The correspondence is justified by the three
properties listed above: intrinsicness, perspectival relativity, and
operator-level inseparability from spatial structure.

Second, the identification of emergent geometry with physical structure
is restricted to the level of spacetime geometry. It does not claim to
capture all of physics. In particular, it does not address the relation
between consciousness and specific material substrates such as neural
tissue. The claim is that at the level of structural categories relevant
to the hard problem, emergent geometry is the right proxy for the
objective-physical side. Whether this extends to more fine-grained
physical descriptions is an open question that lies beyond the scope of
the present paper.

\section{Algebraic Framework and
Axioms}\label{algebraic-framework-and-axioms}

We now state the minimal algebraic framework in which the failure of
independent deformability can be proved. The point of this section is
not to give a general introduction to Type III\(_1\) factors or to
algebraic quantum field theory. It is to isolate the exact package of
assumptions needed for the Phase 1 theorem and to indicate, in each
case, why that assumption belongs in the present hard-problem-led
setting.

\subsection{Ambient Datum and Admissible
Object}\label{ambient-datum-and-admissible-object}

We work in the standard form of an ambient von Neumann algebra \[
(\mathcal{H}_U,\mathcal{A}_U,\Omega),
\] where \(\Omega\) is cyclic and separating for \(\mathcal{A}_U\).
This is the minimal regularity needed for Tomita-Takesaki theory and is
satisfied in all standard AQFT models~\cite{Haag1996}.
Inside \(\mathcal{A}_U\) we select an accessible von Neumann subalgebra
\[
\mathcal{A}_c\subset \mathcal{A}_U,
\] and define the reference state \[
\omega(X):=\langle \Omega, X\Omega\rangle.
\]

The pair \[
(\mathcal{A}_c,\omega)
\] is the basic object of the theory. In the intended applications,
\(\mathcal{A}_c\) is of Type III\(_1\). This is not an arbitrary
technical preference. Type III\(_1\) is the natural habitat of local
quantum field theoretic algebras~\cite{Haag1996}, and it is precisely the setting in
which modular theory is rich enough to support the present route from
algebraic accessibility data to emergent geometry.

The ambient algebra \(\mathcal{A}_U\) should be thought of as the
unreduced field of possible observables, while \(\mathcal{A}_c\) is the
algebra actually stabilized or selected by the accessibility criterion.
The question of the paper is then asked not about the entire ambient
structure at once, but about what can be extracted from the selected
pair \((\mathcal{A}_c,\omega)\).

\subsection{Inherited Accessibility
Conditions}\label{inherited-accessibility-conditions}

The first two assumptions are inherited from the HAFF
setting~\cite{Liu2026PaperD}.

\textbf{(U1) Faithfulness.}
The restriction of the reference state to the accessible algebra is
faithful: \[
\omega|_{\mathcal{A}_c}\ \text{is faithful}.
\] Equivalently, \(\Omega\) is cyclic and separating for
\(\mathcal{A}_c\).

This is the minimal condition that allows modular theory to be applied
to \((\mathcal{A}_c,\omega)\). In the present philosophical translation,
\((U1)\) says that the allegedly ``first-person'' sector is not attached
to a degenerate algebraic viewpoint. The accessible algebra must support
a genuine intrinsic modular structure.

\textbf{(U3) Maximality.}
The accessible algebra is maximal among von Neumann subalgebras that are
faithful and stable under the decoherence semigroup: \[
\mathcal{E}_t(\mathcal{A}_c)\subset \mathcal{A}_c,
\qquad
\forall t\ge 0,
\] and no strictly larger von Neumann subalgebra satisfying \((U1)\) and
this forward stability property exists.

Here \(\{\mathcal{E}_t\}_{t\ge 0}\) is the decoherence semigroup acting
on the ambient algebra: a strongly continuous one-parameter semigroup of
normal completely positive unital maps on \(\mathcal{A}_U\) satisfying
the Markov property \(\mathcal{E}_0=\mathrm{id}\) and
\(\mathcal{E}_{s+t}=\mathcal{E}_s\circ\mathcal{E}_t\). Its physical
role is to model the progressive loss of coherence between the
accessible algebra and its environment; the fixed-point algebra of the
semigroup is \(\mathcal{A}_c\) itself. (For the full construction and
motivation, see~\cite{Liu2026PaperD}.)

The role of \((U3)\) is to prevent an arbitrary shrinking of the
accessible algebra. It says that \(\mathcal{A}_c\) is not merely one
convenient observational slice among many. It is maximal with respect to
the stability criterion that defines accessibility in the first place.
This matters philosophically because the inseparability theorem is meant
to concern the selected structural viewpoint itself, not an artificially
weakened fragment of it.

Throughout this paper, we work within the admissible class selected by
\((U3)\). The Phase 1 proofs below do not use maximality directly; its
role is to ensure that the algebraic datum is canonical rather than
artificially restricted.

\subsection{The New Dynamical Axiom}\label{the-new-dynamical-axiom}

The crucial new input is the scaling relation between decoherence and
modular evolution.

\textbf{(C) Modular-decoherence scaling covariance.}
For all \(s\in\mathbb{R}\) and \(t\ge 0\), \[
\mathcal{E}_t\circ \mathrm{Ad}\,\Delta_{\mathcal{A}_c}^{is}
=
\mathrm{Ad}\,\Delta_{\mathcal{A}_c}^{is}\circ \mathcal{E}_{e^{2\pi s}t}.
\]

Here \(\mathrm{Ad}\,\Delta_{\mathcal{A}_c}^{is}\) denotes the outer
conjugation action of the modular unitary on the ambient operator
algebra.

This is the only genuinely new dynamical axiom introduced in Phase 1.
Its content is that modular evolution does not merely preserve the
accessible algebra tautologically. It rescales the effective decoherence
depth. In the language of the hard problem, \((C)\) is the first place
where the object-side and act-side structures are explicitly tied
together. It is therefore the real replacement for the earlier
non-scrambling intuition.

\begin{remark}[Physical plausibility of \((C)\)]
Axiom \((C)\) holds whenever the modular flow of the accessible algebra
acts as a dilation on the decoherence time parameter. This is the case
in the Bisognano-Wichmann setting, where modular flow is a Lorentz
boost and decoherence depth scales with the boost parameter, and more
generally in any Rindler-type wedge algebra. The axiom is a nontrivial
dynamical input: it ties the intrinsic modular time of the accessible
algebra to the extrinsic decoherence structure of the ambient algebra.
Its domain of validity beyond the HAFF setting~\cite{Liu2026PaperD} is
an open question.
\end{remark}

\subsection{Ambient Shell Realization}\label{ambient-shell-realization}

To turn \((C)\) into a geometric statement, one needs an external family
on which the rescaling can act.

\textbf{(Amb) Ambient shell realization.}
There exists a family of von Neumann algebras \[
\{\mathcal{A}(\tau)\}_{\tau>0}
\] such that:

\begin{enumerate}
\def\labelenumi{\arabic{enumi}.}
\tightlist
\item
  for every \(\tau>0\), \[
  \mathcal{A}_c\subset \mathcal{A}(\tau)\subset \mathcal{A}_U;
  \]
\item
  if \(\tau_1<\tau_2\), then \[
  \mathcal{A}(\tau_1)\supset \mathcal{A}(\tau_2);
  \]
\item
  the accessible algebra is the core limit: \[
  \mathcal{A}_c=\bigcap_{\tau>0}\mathcal{A}(\tau);
  \]
\item
  the family \(\mathcal{A}(\tau)\) is an ambient realization of the
  decoherence shells, so that the scaling computation induced by \((C)\)
  legitimately reads the parameter \(\tau\) as a decoherence window.
\end{enumerate}

The point of \((Amb)\) is regularity, not extra dynamics. It provides a
shell family external to \(\mathcal{A}_c\) but still anchored in the
ambient algebra. Without such a realization, the scaling relation of
\((C)\) would remain formal. With it, modular flow acquires a genuine
action on a stratified family of observational stability levels. This is
the step that allows ``information stability'' to be re-read as
``geometric organization.''

It is important to note what \((Amb)\) is not. It is not an arbitrary
filtration appended for convenience. It is the rigorous formalization of
the shell notion actually used downstream. In particular, outer
standardness will later be derived from \((Amb)\) rather than postulated
separately.

\subsection{Existential Conormality}\label{existential-conormality}

The final input required for the Phase 1 conformal reconstruction is a
second-direction existence condition.

\textbf{(CN\(_\exists\)) Existential conormality.}
There exists some \(\mu_0>0\) such that, for \[
\mathcal{C}_{\mu_0}:=\mathcal{B}_{\mu_0}'\cap \mathcal{A}_c,
\] the vector \(\Omega\) is cyclic.

The algebras \(\mathcal{B}_{\mu}\) will be constructed from \((C)\) and
\((Amb)\) in Section 4; the present point is only that one cyclicity
assumption at one scale is enough to close the second direction and
trigger the GLW reconstruction.

Philosophically, \((CN_\exists)\) is the minimal extra input needed to
pass from a single half-sided modular direction to a full conformal
skeleton. It should therefore be read neither as an arbitrary technical
patch nor as a hidden second dynamics. It is an existence-type condition
ensuring that the accessible algebra contains enough internal relational
richness to support the second lightlike direction.

\subsection{What Is Explicitly Not
Assumed}\label{what-is-explicitly-not-assumed}

One older condition is intentionally removed.

\textbf{(U2) is deleted.}
The statement that the modular group preserves the accessible algebra,
\[
\sigma_t^{\mathcal{A}_c}(\mathcal{A}_c)=\mathcal{A}_c,
\] is automatic from Tomita-Takesaki theory~\cite{Takesaki2003} and therefore carries no
genuine selection content.

This deletion matters conceptually. The paper does not rest on a
tautological ``modular stability'' assumption. The real content enters
through \((C)\), where modular evolution acts outwardly on the shell
family rather than merely inwardly on \(\mathcal{A}_c\) itself.

It is also worth noting that the stronger exhaustion condition \[
\mathrm{OE}: \quad W^*\!\Big(\bigcup_{\tau>0}\mathcal{A}(\tau)\Big)=\mathcal{B}(\mathcal{H}_U)
\] is not part of the minimal Phase 1 package. It is only relevant for
stronger scalar-core statements that are not needed for the present
conformal-level theorem.

\subsection{The Minimal Phase 1
Package}\label{the-minimal-phase-1-package}

We may now summarize the standing assumptions. Throughout the paper we
work within the admissible class selected by \((U3)\). The
proof-relevant Phase 1 package is: \[
(U1)+(C)+(Amb)+(CN_\exists).
\]

The logic of this package is straightforward.

\begin{itemize}
\tightlist
\item
  \((U1)\) makes modular theory available on the accessible algebra.
\item
  \((C)\) ties modular evolution to decoherence depth.
\item
  \((Amb)\) gives that depth a concrete shell realization inside the
  ambient algebra.
\item
  \((CN_\exists)\) provides the minimal extra cyclicity needed to
  reconstruct the second direction and hence the conformal skeleton.
\end{itemize}

The remainder of Phase 1 consists in proving that this proof-relevant
package is
already enough to derive a standard half-sided modular inclusion, a
canonical GLW Möbius net, and therefore a conformal-level inseparability
theorem. That is the task of the next section.

\section{From Axioms to Conformal Skeleton (Phase
1)}\label{from-axioms-to-conformal-skeleton-phase-1}

\subsection{From Axiom C to Standard Half-Sided Modular
Inclusion}\label{from-axiom-c-to-standard-half-sided-modular-inclusion}

The first structural point of the paper is that the half-sided modular
inclusion needed by the Wiesbrock-GLW machinery is not postulated. It is
extracted from the HAFF-native dynamics once the modular-decoherence
covariance axiom is combined with an ambient realization of the
decoherence shells. In particular, the present route differs sharply
from the earlier formulation in which HSMI appeared as an external
ansatz. Here it is the output of the Phase 1 assumptions.

Throughout this section we assume \((U1)\), \((C)\), and \((Amb)\), and
we work within the admissible class selected by \((U3)\). We write \[
\sigma_t := \sigma_t^{\mathcal{A}_c}
\] for the modular automorphism group of \((\mathcal{A}_c,\Omega)\) and
\[
J := J_{\mathcal{A}_c}
\] for the associated modular conjugation. By \((Amb)\), we are given a
decreasing family of ambient shell algebras \[
\mathcal{A}_c \subset \mathcal{A}(\tau) \subset \mathcal{A}_U,
\qquad
\tau>0,
\] with \[
\mathcal{A}_c=\bigcap_{\tau>0}\mathcal{A}(\tau),
\] and whose parameter \(\tau\) is interpreted as a decoherence window.

\subsection{Outer Scaling from Modular-Decoherence
Covariance}\label{outer-scaling-from-modular-decoherence-covariance}

The real content of axiom \((C)\) is that the modular flow of the
accessible algebra does not merely preserve \(\mathcal{A}_c\)
tautologically. Rather, it rescales the ambient shell parameter. This is
the first point at which the older \((U2)\)-style formulation must be
replaced.

\begin{proposition}[Outer scaling law]\label{prop:4.1}
Under \((U1)\), \((C)\), and \((Amb)\), \[
\sigma_s(\mathcal{A}(\tau))=\mathcal{A}(e^{-2\pi s}\tau),
\qquad
\forall s\in\mathbb{R},\ \tau>0.
\]
\end{proposition}

\begin{proof}
By the shell interpretation built into \((Amb)\), membership in
\(\mathcal{A}(\tau)\) is determined by stability on the time window of
size \(\tau\). For any \(O\in\mathcal{A}(\tau)\), axiom \((C)\) gives \[
\mathcal{E}_t\!\big(\sigma_s(O)\big)
=
\sigma_s\!\big(\mathcal{E}_{e^{2\pi s}t}(O)\big),
\] hence the same stability estimate that characterizes \(O\) on the
window \(e^{2\pi s}t\le \tau\) characterizes \(\sigma_s(O)\) on the
window \(t\le e^{-2\pi s}\tau\). Therefore
\(\sigma_s(O)\in\mathcal{A}(e^{-2\pi s}\tau)\). Applying the same
argument to \(\sigma_{-s}\) yields the reverse inclusion.
\end{proof}

This result is the correct algebraic replacement for the old
non-scrambling intuition: modular time rescales decoherence depth. It is
already an inseparability statement in weak form, because the shell
parameter and the modular parameter are no longer independently
adjustable.

\subsection{Outer Standardness}\label{outer-standardness}

The next observation is simple but essential: the ambient shells are
standard for the same vector \(\Omega\). This is where the ambient
realization is used as a regularity hypothesis rather than as extra
dynamics.

\begin{lemma}[Outer standardness]\label{lem:4.2}
For every \(\tau>0\), the vector \(\Omega\) is cyclic and separating for
\(\mathcal{A}(\tau)\).
\end{lemma}

\begin{proof}
Since \(\mathcal{A}_c\subset \mathcal{A}(\tau)\) and
\(\Omega\) is cyclic for \(\mathcal{A}_c\) by \((U1)\), it is cyclic for
\(\mathcal{A}(\tau)\). Since \(\mathcal{A}(\tau)\subset \mathcal{A}_U\)
and \(\Omega\) is separating for \(\mathcal{A}_U\) by the ambient datum
(Section 3.1), separation passes to the subalgebra
\(\mathcal{A}(\tau)\).
\end{proof}

This standardness will allow us to transport the outer shell family into
an internal family inside \(\mathcal{A}_c\) by Tomita-Takesaki duality.

\subsection{\texorpdfstring{The \(J\)-Dual Internal
Filtration}{The J-Dual Internal Filtration}}\label{the-j-dual-internal-filtration}

The key construction is to turn the outer shell family into a family of
subalgebras of the accessible algebra itself. This is the point at which
the earlier interiorization axiom becomes unnecessary.

For \(\mu>0\), define \[
\mathcal{B}_\mu:=J\,\mathcal{A}(1/\mu)'\,J.
\]

The next proposition records the properties of this family needed for
the HSMI step.

\begin{proposition}[\(J\)-dual internal filtration]\label{prop:4.3}
The family \(\{\mathcal{B}_\mu\}_{\mu>0}\) satisfies:

\begin{enumerate}
\def\labelenumi{\arabic{enumi}.}
\tightlist
\item
  \(\mathcal{B}_\mu\subset \mathcal{A}_c\) for all \(\mu>0\);
\item
  if \(\mu_1>\mu_2\), then
  \(\mathcal{B}_{\mu_1}\subset \mathcal{B}_{\mu_2}\);
\item
  \(\Omega\) is cyclic and separating for each \(\mathcal{B}_\mu\);
\item
  the modular group rescales the family according to \[
  \sigma_s(\mathcal{B}_\mu)=\mathcal{B}_{e^{-2\pi s}\mu};
  \]
\item
  the family generates the accessible algebra: \[
  W^*\!\Big(\bigcup_{\mu>0}\mathcal{B}_\mu\Big)=\mathcal{A}_c.
  \]
\end{enumerate}
\end{proposition}

\begin{proof}
The inclusion \(\mathcal{A}_c\subset \mathcal{A}(1/\mu)\) implies \[
\mathcal{A}(1/\mu)'\subset \mathcal{A}_c',
\] hence \[
\mathcal{B}_\mu
=
J\,\mathcal{A}(1/\mu)'\,J
\subset
J\,\mathcal{A}_c'\,J
=
\mathcal{A}_c,
\] which proves (1).

For (2), if \(\mu_1>\mu_2\), then \(1/\mu_1<1/\mu_2\), so the
monotonicity of the outer shell family gives \[
\mathcal{A}(1/\mu_1)\supset \mathcal{A}(1/\mu_2),
\] hence \[
\mathcal{A}(1/\mu_1)'\subset \mathcal{A}(1/\mu_2)'
\] and therefore \[
\mathcal{B}_{\mu_1}\subset \mathcal{B}_{\mu_2}.
\]

For (3), Lemma 4.2 shows that \(\Omega\) is cyclic and separating for
\(\mathcal{A}(1/\mu)\). Therefore \(\Omega\) is also cyclic and
separating for \(\mathcal{A}(1/\mu)'\). Since \(J\Omega=\Omega\) and
\(J\) is antiunitary, the same holds for \[
J\,\mathcal{A}(1/\mu)'\,J=\mathcal{B}_\mu.
\]

For (4), using the Tomita relation \[
\Delta^{is}J=J\Delta^{-is},
\] we compute \[
\sigma_s(\mathcal{B}_\mu)
=
\Delta^{is}J\,\mathcal{A}(1/\mu)'\,J\Delta^{-is}
=
J\Delta^{-is}\mathcal{A}(1/\mu)'\Delta^{is}J
=
J\,\sigma_{-s}\!\big(\mathcal{A}(1/\mu)\big)'\,J.
\] Now Proposition 4.1 gives \[
\sigma_{-s}\!\big(\mathcal{A}(1/\mu)\big)
=
\mathcal{A}(e^{2\pi s}/\mu),
\] so \[
\sigma_s(\mathcal{B}_\mu)
=
J\,\mathcal{A}(e^{2\pi s}/\mu)'\,J
=
\mathcal{B}_{e^{-2\pi s}\mu}.
\]

Finally, (5) follows from the core limit in \((Amb)\): \[
\mathcal{A}_c=\bigcap_{\tau>0}\mathcal{A}(\tau)
\quad\Longrightarrow\quad
\mathcal{A}_c'=W^*\!\Big(\bigcup_{\tau>0}\mathcal{A}(\tau)'\Big).
\] Applying \(J(\cdot)J\) yields \[
\mathcal{A}_c
=
J\mathcal{A}_c'J
=
W^*\!\Big(\bigcup_{\tau>0}J\,\mathcal{A}(\tau)'\,J\Big)
=
W^*\!\Big(\bigcup_{\mu>0}\mathcal{B}_\mu\Big).
\]
\end{proof}

This construction is the first genuinely new structural step of the
Phase 1 chain. The internal family is not imposed from outside. It is
obtained functorially from the shell realization and modular data.

\subsection{Half-Sided Modular Inclusion as a
Theorem}\label{half-sided-modular-inclusion-as-a-theorem}

We may now isolate the fixed inclusion that drives the Wiesbrock
machinery.

\begin{theorem}[HSMI from modular-decoherence covariance]\label{thm:4.4}
For any fixed \(\mu_0>0\), the inclusion \[
(\mathcal{B}_{\mu_0}\subset \mathcal{A}_c,\Omega)
\] is a half-sided modular inclusion, namely \[
\sigma_t(\mathcal{B}_{\mu_0})\subset \mathcal{B}_{\mu_0},
\qquad
\forall t\le 0.
\]
\end{theorem}

\begin{proof}
By Proposition 4.3(3), \(\Omega\) is cyclic and separating for both
\(\mathcal{B}_{\mu_0}\) and \(\mathcal{A}_c\). By Proposition 4.3(4), \[
\sigma_t(\mathcal{B}_{\mu_0})
=
\mathcal{B}_{e^{-2\pi t}\mu_0}.
\] If \(t\le 0\), then \(e^{-2\pi t}\ge 1\), hence
\(e^{-2\pi t}\mu_0\ge \mu_0\). By Proposition 4.3(2), \[
\mathcal{B}_{e^{-2\pi t}\mu_0}\subset \mathcal{B}_{\mu_0}.
\] Therefore \[
\sigma_t(\mathcal{B}_{\mu_0})\subset \mathcal{B}_{\mu_0},
\qquad t\le 0.
\]
\end{proof}

The important point is conceptual as much as technical: in the present
framework, HSMI is not a geometric input imported from AQFT\@. It is the
algebraic consequence of a scaling law that ties modular flow to
decoherence depth.

\subsection{The Wiesbrock-Borchers Positive
Generator}\label{the-wiesbrock-borchers-positive-generator}

Once the fixed inclusion is available, the standard Wiesbrock-Borchers
conclusion follows.

\begin{corollary}[Positive translation generator for the fixed
inclusion]\label{cor:4.5}
For each fixed \(\mu_0>0\), there exists a unique positive self-adjoint
operator \[
P_{\mu_0}\ge 0
\] such that, with \[
U_{\mu_0}(a):=e^{-iaP_{\mu_0}},
\qquad
a\ge 0,
\] one has

\begin{enumerate}
\def\labelenumi{\arabic{enumi}.}
\tightlist
\item
  \(U_{\mu_0}(a)\Omega=\Omega\) for all \(a\ge 0\);
\item
  the Borchers covariance relation \[
  \Delta_{\mathcal{A}_c}^{is}U_{\mu_0}(a)\Delta_{\mathcal{A}_c}^{-is}
  =
  U_{\mu_0}(e^{-2\pi s}a)
  \] holds for all \(s\in\mathbb{R}\) and \(a\ge 0\).
\end{enumerate}
\end{corollary}

\begin{proof}
This is the standard Wiesbrock theorem for a fixed half-sided modular
inclusion~\cite{Wiesbrock1993}, together with Borchers covariance~\cite{Borchers1992}.
\end{proof}

At this stage we do not claim a pointwise identification between the
entire family \(\{\mathcal{B}_\mu\}_{\mu>0}\) and the Wiesbrock orbit
generated by \(P_{\mu_0}\). Such an identification requires an
additional normalization or uniqueness argument and is not needed for
the Phase 1 main chain.

This completes the first half of Phase 1. We have shown that the
accessible algebra canonically carries the half-sided modular structure
needed to enter the Möbius reconstruction theorem.

\subsection{GLW Conformal Skeleton and Weak
Inseparability}\label{glw-conformal-skeleton-and-weak-inseparability}

The second half of Phase 1 upgrades the fixed half-sided modular
inclusion to a full conformal skeleton. The only additional input is an
existential conormality condition. Once this is imposed, the
Guido-Longo-Wiesbrock reconstruction theorem produces a canonical
strongly additive local Möbius net on \(S^1\), and with it the
conformal-level version of inseparability.

Throughout this section we assume, in addition to the hypotheses of
Section 4.1, the conormality condition \((CN_\exists)\): there exists
\(\mu_0>0\) such that the vector \(\Omega\) is cyclic for \[
\mathcal{C}_{\mu_0}:=\mathcal{B}_{\mu_0}'\cap \mathcal{A}_c.
\]

\subsection{The Second Direction from Existential
Conormality}\label{the-second-direction-from-existential-conormality}

The role of \((CN_\exists)\) is very precise. It is not used to
construct the half-sided inclusion itself. That has already been
achieved in Section 4.1. Rather, it upgrades the fixed inclusion to a
standard one in the sense required by the GLW theorem.

For \(\mu>0\), define \[
\mathcal{C}_\mu:=\mathcal{B}_\mu'\cap \mathcal{A}_c.
\]

\begin{proposition}[The second direction]\label{prop:4.6}
Under \((U1)\), \((C)\), \((Amb)\), and \((CN_\exists)\), the
family \(\{\mathcal{C}_\mu\}_{\mu>0}\) satisfies:

\begin{enumerate}
\def\labelenumi{\arabic{enumi}.}
\tightlist
\item
  \(\mathcal{C}_\mu\subset \mathcal{A}_c\) and \(\Omega\) is separating
  for \(\mathcal{C}_\mu\);
\item
  the modular group rescales the family by \[
  \sigma_t(\mathcal{C}_\mu)=\mathcal{C}_{e^{-2\pi t}\mu};
  \]
\item
  for all \(t\ge 0\), \[
  \sigma_t(\mathcal{C}_\mu)\subset \mathcal{C}_\mu;
  \]
\item
  if \(\Omega\) is cyclic for one \(\mathcal{C}_{\mu_0}\), then by
  modular propagation it is cyclic for all \(\mathcal{C}_\mu\).
\end{enumerate}

In particular, the fixed inclusion \[
(\mathcal{B}_{\mu_0}\subset \mathcal{A}_c,\Omega)
\] is standard.
\end{proposition}

\begin{proof}
The inclusion \(\mathcal{C}_\mu\subset \mathcal{A}_c\) is immediate from
the definition. Separation of \(\Omega\) follows from \((U1)\), since
\(\omega|_{\mathcal{A}_c}\) is faithful and faithfulness passes to
subalgebras.

For modular covariance, \[
\sigma_t(\mathcal{C}_\mu)
=
\sigma_t(\mathcal{B}_\mu'\cap \mathcal{A}_c)
=
\sigma_t(\mathcal{B}_\mu)'\cap \sigma_t(\mathcal{A}_c)
=
\mathcal{B}_{e^{-2\pi t}\mu}'\cap \mathcal{A}_c
=
\mathcal{C}_{e^{-2\pi t}\mu},
\] where we used Proposition 4.3(4) and the tautological invariance of
\(\mathcal{A}_c\) under its own modular group.

If \(t\ge 0\), then \(e^{-2\pi t}\le 1\), so the monotonicity of the
\(\mathcal{B}_\mu\) family implies the reverse monotonicity of the
relative commutants: \[
\mathcal{C}_{e^{-2\pi t}\mu}\subset \mathcal{C}_\mu.
\] Thus \[
\sigma_t(\mathcal{C}_\mu)\subset \mathcal{C}_\mu.
\]

Finally, \((CN_\exists)\) gives cyclicity for one
\(\mathcal{C}_{\mu_0}\), and the modular covariance formula propagates
this cyclicity along the entire orbit. Since Section 4.1 already gives
cyclicity and separation for \(\mathcal{B}_{\mu_0}\) and
\(\mathcal{A}_c\), the inclusion is standard.
\end{proof}

This is the exact point at which the second direction closes. No
stronger global conormality assumption is needed for the fixed inclusion
that feeds the GLW theorem.

\subsection{The GLW Reconstruction
Theorem}\label{the-glw-reconstruction-theorem}

Once standardness is available, the conformal skeleton appears by a
known AQFT theorem. The novelty in the present paper is not the GLW
theorem itself, but the fact that its input is derived from the
HAFF-native algebraic package.

\begin{theorem}[Conformal skeleton via Guido-Longo-Wiesbrock]\label{thm:4.7}
Fix \(\mu_0>0\) as in \((CN_\exists)\). Then the standard half-sided
modular inclusion \[
(\mathcal{B}_{\mu_0}\subset \mathcal{A}_c,\Omega)
\] determines a unique strongly additive local Möbius covariant net \[
I\longmapsto \mathcal{A}_{\mathrm{GLW}}(I)
\] on \(S^1\) such that \[
\mathcal{A}_{\mathrm{GLW}}(0,+\infty)=\mathcal{A}_c,
\qquad
\mathcal{A}_{\mathrm{GLW}}(1,+\infty)=\mathcal{B}_{\mu_0},
\] and \(\Omega\) is the vacuum vector for the corresponding
positive-energy representation of \(\mathrm{PSL}(2,\mathbb{R})\).
\end{theorem}

\begin{proof}
This is precisely the Guido-Longo-Wiesbrock reconstruction theorem for
standard half-sided modular inclusions~\cite{GLW1998}.
\end{proof}

Two clarifications are important here.

First, at this stage no separate relative-position hypothesis is needed.
Standard HSMI is enough to reconstruct the Möbius net. The previously
suspected relative-position bottleneck concerns later identification
questions, not the existence of the conformal skeleton itself.

Second, the net produced here is an abstract but canonical conformal
object attached to the fixed inclusion. The remaining question of
identifying this GLW net pointwise with the pre-given shell filtrations
or later cross-ratio families belongs to the next phase of the program.

\subsection{The Conformal Skeleton as the Output of Modular
Data}\label{the-conformal-skeleton-as-the-output-of-modular-data}

Theorem 4.7 already yields the first inseparability statement of the
paper. The data entering the standard inclusion are modular data; the
output is a conformal geometry in the minimal sense appropriate to one
spatial and one temporal dimension, namely a causal/conformal skeleton.

To make this precise, let \[
\mathrm{Mod}_{\mathrm{conf}}(\mathcal{A}_c,\omega)
\] denote the conformal modular sector determined in Phase 1: the
modular group of \((\mathcal{A}_c,\Omega)\), the fixed-inclusion
positive generator \(P_{\mu_0}\), and the corresponding Möbius
representation recovered from the standard HSMI\@. Let \[
\mathrm{Conf}(\mathcal{A}_c,\omega)
\] denote the conformal skeleton, equivalently the GLW net together with
its Möbius vacuum representation.

\begin{theorem}[Weak Inseparability]\label{thm:4.8}
Within the admissible class selected by \((U3)\), and under \((U1)\),
\((C)\), \((Amb)\), and \((CN_\exists)\), the conformal skeleton and the
conformal modular sector cannot be independently varied. More precisely:

\begin{enumerate}
\def\labelenumi{\arabic{enumi}.}
\tightlist
\item
  the standard HSMI data determine a unique conformal skeleton;
\item
  the conformal skeleton, once normalized by the half-line
  identification above, determines the conformal modular sector;
\item
  therefore any admissible deformation that preserves one of these two
  structures must preserve the other.
\end{enumerate}
\end{theorem}

\begin{proof}
Item (1) is exactly Theorem 4.7.

For item (2), the Möbius net comes equipped with a positive-energy
representation of \(\mathrm{PSL}(2,\mathbb{R})\). Once the half-line
normalization \[
\mathcal{A}_{\mathrm{GLW}}(0,+\infty)=\mathcal{A}_c,
\qquad
\mathcal{A}_{\mathrm{GLW}}(1,+\infty)=\mathcal{B}_{\mu_0}
\] is fixed, the corresponding dilation and translation subgroups
recover the modular action and the positive translation generator
associated with the inclusion. In other words, the conformal skeleton
rigidly determines the conformal modular sector.

Item (3) is immediate. If an admissible deformation preserves the
conformal modular sector, item (1) forces the conformal skeleton to
remain fixed. Conversely, if it preserves the conformal skeleton, item
(2) forces the modular sector to remain fixed. Hence the two structures
are not independently deformable.
\end{proof}

This result is weak only in the sense that it stops at the conformal
level. It does not yet recover metric scale. That is the purpose of Step
I and Step II\@. But it already establishes the central conceptual point
of Phase 1:

\begin{quote}
causal/conformal geometry and modular time are not two separately
specified ingredients;\\
they are two readings of the same standard half-sided modular data.
\end{quote}

\subsection{What Phase 1 Has Actually
Closed}\label{what-phase-1-has-actually-closed}

It is useful to isolate the precise mathematical output of the section
before moving on to metric scale.

\begin{enumerate}
\def\labelenumi{\arabic{enumi}.}
\tightlist
\item
  HSMI is no longer an assumption. It is a theorem derived from \((C)\)
  and \((Amb)\).
\item
  The second direction is no longer a global mystery. It is isolated
  into the existence-type conormality input \((CN_\exists)\).
\item
  The conformal skeleton is no longer conjectural. It is reconstructed
  canonically by GLW from the fixed standard inclusion.
\item
  The first inseparability theorem is already in hand: conformal
  geometry and conformal modular data are mutually rigid.
\end{enumerate}

What remains open after Phase 1 is not the existence of the conformal
skeleton, but the recovery of metric scale from local modular data. That
is why the next phase turns to the cross-ratio family and to the
canonical local modular scale datum \[
\{K_b\}_{b\in(0,1)}.
\]

\section{From Conformal Skeleton to Full Geometry (Phase
2)}\label{from-conformal-skeleton-to-full-geometry-phase-2}

Phase 1 established that conformal geometry and conformal modular data
are inseparable. But the conformal skeleton determines only the causal
structure, the light cones. It does not determine distances. A conformal
class of Lorentzian metrics contains infinitely many metrics related by
Weyl rescaling. To reach the full inseparability statement, we must show
that the metric scale is also locked to modular data.

This is the task of Phase 2. The strategy has three steps: identify the
correct local geometric modulus, extract a canonical modular scale datum
from it, and show that this datum determines the metric scale via a Type
III-compatible quantum Fisher information metric.

\subsection{Why Single-Interval Deformations
Fail}\label{why-single-interval-deformations-fail}

The first instinct is to deform a single interval within the GLW net and
track how the modular Hamiltonian changes. This fails for a precise
reason.

On \(S^1\), the group \(\mathrm{PSL}(2,\mathbb{R})\) acts transitively
on the space of intervals. Any single-interval deformation, such as
moving one endpoint while fixing the other, is a Möbius transformation.
Since the vacuum state \(\omega\) is Möbius-invariant, pulling back the
restricted state along such a deformation produces a constant family.
The quantum Fisher information metric of a constant family is zero.

The lesson is that single-interval endpoint motion is a gauge orbit, not
a geometric deformation. To obtain a nontrivial modulus, one needs at
least four boundary points.

\subsection{The Cross-Ratio as the Minimal Non-Möbius
Modulus}\label{the-cross-ratio-as-the-minimal-non-muxf6bius-modulus}

Consider a nested double-interval configuration within the GLW net: \[
\mathcal{M}_{\mathrm{in}}(b) \subset \mathcal{M}_{\mathrm{out}},
\qquad b \in (0,1).
\] In the standard half-line gauge, this means \[
\mathcal{M}_{\mathrm{out}} = \mathcal{A}_{\mathrm{GLW}}(0,\infty),
\qquad
\mathcal{M}_{\mathrm{in}}(b) = \mathcal{A}_{\mathrm{GLW}}(b,1).
\] The four boundary points are \(0, b, 1, \infty\). The Möbius group
has three real parameters, so after gauge-fixing, exactly one real
modulus survives: the cross-ratio \[
\chi = b.
\]

This is the minimal non-Möbius geometric degree of freedom available in
the GLW net. Varying \(b\) genuinely changes the relative position of
the inner and outer algebras in a way that no Möbius transformation can
undo.

\subsection{The Canonical Modular Scale
Datum}\label{the-canonical-modular-scale-datum}

For each \(b \in (0,1)\), the inclusion
\(\mathcal{M}_{\mathrm{in}}(b) \subset \mathcal{M}_{\mathrm{out}}\)
determines a modular operator
\(\Delta_{\mathcal{M}_{\mathrm{in}}(b),\Omega}\) via the Tomita-Takesaki
theorem applied to \(\mathcal{M}_{\mathrm{in}}(b)\) with respect to the
vacuum vector \(\Omega\).

\begin{definition}[Canonical modular scale datum]\label{def:5.1}
\[
K_b := -\ln \Delta_{\mathcal{M}_{\mathrm{in}}(b),\Omega},
\qquad b \in (0,1).
\]
\end{definition}

The family \(\{K_b\}_{b \in (0,1)}\) is the canonical local modular
scale datum. It is determined entirely by the GLW net and the vacuum
state. No additional choices are involved.

\begin{proposition}[Cross-ratio is the minimal non-M\"obius
modulus]\label{prop:5.2}
A nested double-interval configuration in the GLW net has
exactly one Möbius-invariant modulus, the cross-ratio \(\chi = b\). In
particular, the family \(\{K_b\}\) captures the minimal local modular
scale information that is not already contained in the conformal
skeleton.
\end{proposition}

This is the theorem-level output of Step I. It requires no additional
axioms beyond Phase 1.

\subsection{The Chart Axiom and the Scale
Shadow}\label{the-chart-axiom-and-the-scale-shadow}

To convert the modular scale datum into a Riemannian metric on the
cross-ratio parameter space, we need to differentiate the family
\(\{K_b\}\) and evaluate a quadratic form. This requires regularity that
is not automatic from the abstract GLW construction.

\textbf{Axiom (LImp\(^1\)).} For each reference point \(b_0 \in (0,1)\),
there exists a neighborhood \(U_{b_0}\) and a unitary family \[
W_b \in \mathcal{U}(\mathcal{M}_{\mathrm{out}}),
\qquad b \in U_{b_0},
\] satisfying:

\begin{enumerate}
\def\labelenumi{\arabic{enumi}.}
\tightlist
\item
  \(W_{b_0} = \mathbf{1}\);
\item
  \(W_b \mathcal{M}_* W_b^* = \mathcal{M}_{\mathrm{in}}(b)\) for a fixed
  reference algebra \(\mathcal{M}_* = \mathcal{M}_{\mathrm{in}}(b_0)\);
\item
  \(b \mapsto \omega(W_b A W_b^*)\) is \(C^1\) for each
  \(A \in \mathcal{M}_*\);
\item
  the tangent vector \(\dot{\phi}_{b_0}\) at \(b = b_0\) is represented
  by the modular Hamiltonian variation
  \(\dot{K}_{b_0} = \frac{d}{db}|_{b=b_0} K_b\) via the standard Araki
  correspondence.
\end{enumerate}

This axiom packages three regularity conditions into one: the existence
of a local implementer (CR1), \(C^1\) smoothness of the pulled-back
state family (CR2), and tangent compatibility with the modular
Hamiltonian variation (CR3).

\textbf{Axiom (CR-nd) (Non-degeneracy).} The Araki-Hessian quadratic
form \(\mathbf{Q}_{\omega,b}\) is strictly positive on the modular
Hamiltonian variation: \[
\mathbf{Q}_{\omega,b}(\partial_b K_b, \partial_b K_b) > 0,
\qquad \forall\, b \in (0,1).
\]
In the worked examples, this follows from the positivity of the
stress-tensor two-point function. It fails only if the cross-ratio
deformation produces no state-level change, a physically degenerate
situation.

Under these two conditions, the scale shadow is a genuine Riemannian
metric on the cross-ratio parameter space.

\begin{theorem}[Scale shadow, conditional on \textnormal{(LImp\(^1\))} +
\textnormal{(CR-nd)}]\label{thm:5.3}
Under the Phase 1 axioms and the additional conditions
\((LImp^1)\) and \((CR\text{-}nd)\), the canonical modular scale datum
induces a Riemannian scale shadow \[
h_b = \mathbf{Q}_{\omega,b}(\partial_b K_b, \partial_b K_b)\, db^2,
\] where \(\mathbf{Q}_{\omega,b}\) is the positive quadratic form
induced by the Araki relative entropy Hessian~\cite{Araki1976} at the state
\(\omega|_{\mathcal{M}_{\mathrm{in}}(b)}\) via the \((LImp^1)\) chart.
\end{theorem}

\begin{remark}
In the special case where the chart is of exponential
type (CR-exp), this reduces to the familiar variance formula \[
h_b = \mathrm{Var}_\omega(\partial_b K_b)\, db^2.
\] The general formula via \(\mathbf{Q}_{\omega,b}\) does not require
this specialization.
\end{remark}

\subsection{Geometric Datum: Conformal Skeleton + Slice
Scale}\label{assembly-conformal-skeleton-scale-lorentzian-geometry}

The conformal skeleton from Phase 1 determines the causal structure,
equivalently the conformal causal structure. The scale shadow
\(\{h_b\}\) determines the complementary scale information along the
cross-ratio slice family. At the abstract level established in this
paper, the pair \((\mathrm{Conf},\{h_b\})\) is the full geometric datum.
In the worked examples of Section 6, the cross-ratio parameter \(b\)
naturally parametrizes the spatial position of the inner interval
endpoint within the outer interval. The scale shadow \(h_b\) therefore
directly defines a Riemannian metric on the spatial slice, and together
with the conformal causal structure determines a Lorentzian metric.

\begin{corollary}[Full geometric datum, 1+1 dimensions]\label{cor:5.4}
In the 1+1
dimensional setting of the GLW net, the conformal skeleton
\(\mathrm{Conf}(\mathcal{N},\omega)\) together with the scale family
\(\{h_b\}_{b \in (0,1)}\) defines the full geometric datum \[
\mathrm{Geom}_{\mathrm{full}}(\mathcal{N},\omega)
:=
\big(\mathrm{Conf}(\mathcal{N},\omega),\; \{h_b\}_{b \in (0,1)}\big).
\]
In the worked examples of Section 6, the cross-ratio parameter \(b\)
naturally parametrizes the spatial position of the inner interval
endpoint within the outer interval. The scale shadow \(h_b\) therefore
directly defines a Riemannian metric on the spatial slice, and together
with the conformal causal structure determines a Lorentzian metric.
\end{corollary}

\begin{remark}
The step from scale shadow to Lorentzian metric relies on the
model-level identification of the cross-ratio parameter with spatial
position. At the abstract level, Corollary~\ref{cor:5.4} provides a
full geometric datum; the Lorentzian interpretation is verified in the
worked examples but not proved in full generality. The extension to
higher dimensions is an open problem (Section~\ref{open-problems}).
\end{remark}

\subsection{Forward and Bidirectional Full
Inseparability}\label{forward-and-bidirectional-full-inseparability}

With the full geometric datum in hand, we can state the main theorems.

\begin{definition}[Extended modular data]\label{def:5.5}
\[
\mathrm{Mod}_{\mathrm{ext}}(\mathcal{N},\omega)
:=
\big(\mathrm{Mod}_{\mathrm{conf}}(\mathcal{N},\omega),\; \{K_b\}_{b \in (0,1)}\big).
\]
\end{definition}

Two extended modular data
\(\mathrm{Mod}_{\mathrm{ext}}^{(1)}\) and
\(\mathrm{Mod}_{\mathrm{ext}}^{(2)}\) are identified if their conformal
modular sectors agree and their canonical scale families \(\{K_b\}\)
agree up to the trivial equivalences specified in \((SFa)\). All
equalities of \(\mathrm{Mod}_{\mathrm{ext}}\) below are understood in
this quotient sense.

\begin{theorem}[Forward Full Inseparability]\label{thm:5.6}
Under the Phase 1
axioms and \((LImp^1) + (CR\text{-}nd)\), if two admissible objects have
the same extended modular data, then they have the same full geometric
datum:
\[
\mathrm{Mod}_{\mathrm{ext}}^{(1)} = \mathrm{Mod}_{\mathrm{ext}}^{(2)}
\quad \Longrightarrow \quad
\mathrm{Geom}_{\mathrm{full}}^{(1)} = \mathrm{Geom}_{\mathrm{full}}^{(2)}.
\]
\end{theorem}

\begin{proof}
By Phase 1 Weak Inseparability, the conformal modular
data determine the conformal skeleton. By Theorem 5.3, the canonical
modular scale family \(\{K_b\}\) determines the scale shadow
\(\{h_b\}\). By Corollary 5.4, the conformal skeleton and scale shadow
together determine the full geometric datum.
\end{proof}

The forward direction says: you cannot change the full geometric datum
without changing the modular data. But can you change the modular data
without changing the full geometric datum? This requires an additional
faithfulness
condition.

\textbf{Axiom (SFa) (Scale faithfulness).} The canonical scale map
\(\{K_b\} \mapsto \{h_b\}\) is faithful within the admissible class: if
\(h_b^{(1)} = h_b^{(2)}\) for all \(b\), then \(K_b^{(1)} = K_b^{(2)}\)
for all \(b\), up to trivial equivalences that do not affect the scale
map.

\begin{theorem}[Bidirectional Full Inseparability]\label{thm:5.7}
Under the
Phase 1 axioms, \((LImp^1)\), \((CR\text{-}nd)\), and \((SFa)\), \[
\mathrm{Geom}_{\mathrm{full}}^{(1)} = \mathrm{Geom}_{\mathrm{full}}^{(2)}
\quad \Longleftrightarrow \quad
\mathrm{Mod}_{\mathrm{ext}}^{(1)} = \mathrm{Mod}_{\mathrm{ext}}^{(2)}.
\] The full geometric datum and the extended modular data are not independently
deformable.
\end{theorem}

\begin{proof}
The forward direction (\(\Leftarrow\)) is Theorem 5.6.
For the reverse direction (\(\Rightarrow\)): by the definition of
\(\mathrm{Geom}_{\mathrm{full}}\), equality of full geometries
determines both the conformal skeleton and the scale family \(\{h_b\}\).
The conformal skeleton determines the conformal modular data by Phase 1
Weak Inseparability. The scale family determines the canonical modular
scale family \(\{K_b\}\) by \((SFa)\). Together,
\(\mathrm{Mod}_{\mathrm{ext}}^{(1)} = \mathrm{Mod}_{\mathrm{ext}}^{(2)}\).
\end{proof}

This is the central mathematical result of the paper. Under the
identifications defended in Section 2, the mathematical counterpart of
physical structure (the full geometric datum
\(\mathrm{Geom}_{\mathrm{full}}\)) and the mathematical counterpart of
first-person time (extended modular data) cannot be independently
specified. The independent-describability presupposition of the hard
problem fails at the full scale-sensitive level.

\section{Abstract Sufficient Conditions for \texorpdfstring{$(LImp^1)$}{(LImp1)}}\label{sufficient-conditions-limp}

The Phase~2 theorems are conditional on the chart axiom $(LImp^1)$, the non-degeneracy condition $(CR\text{-}nd)$, and the faithfulness condition $(SFa)$. Of these three, $(LImp^1)$ is the most structurally demanding: it requires a local unitary implementer family with $C^1$ regularity and tangent compatibility. This section identifies two abstract routes for discharging $(LImp^1)$: a full discharge for Diff-covariant nets, and a decomposition of the remaining non-Diff burden into transportability, inner implementability, and regularity.

The conditions $(CR\text{-}nd)$ and $(SFa)$ are not addressed here. They appear to depend on dynamical and spectral properties of the specific model (e.g., the positivity of the stress-tensor two-point function for $(CR\text{-}nd)$, and the injectivity of the variance map for $(SFa)$), and are better treated as model-level verifications rather than consequences of structural regularity.

\subsection{Decomposition of \texorpdfstring{$(LImp^1)$}{(LImp1)}}\label{decomposition-limp}

It is useful to isolate the three independent requirements packaged into $(LImp^1)$:

\begin{enumerate}
\item \textbf{Local transport.} For each $b$ near a reference point $b_0$, there exists an isomorphism $\alpha_b\colon \mathcal{M}_* \to \mathcal{M}_{\mathrm{in}}(b)$ that extends to an automorphism of $\mathcal{M}_{\mathrm{out}}$.
\item \textbf{Unitary implementability.} The automorphism $\alpha_b$ is inner in $\mathcal{M}_{\mathrm{out}}$: there exists $W_b \in \mathcal{U}(\mathcal{M}_{\mathrm{out}})$ with $\alpha_b = \mathrm{Ad}(W_b)$.
\item \textbf{$C^1$ regularity and tangent compatibility.} The map $b \mapsto \omega(W_b A W_b^*)$ is $C^1$ for each $A \in \mathcal{M}_*$, and the tangent at $b_0$ is represented by the modular Hamiltonian variation $\dot{K}_{b_0}$ via the Araki correspondence.
\end{enumerate}

Different sufficient conditions address different subsets of these requirements.

\subsection{Why Split Property Alone Is Not Sufficient}\label{why-split-alone-insufficient}

The split property for the inclusion $\mathcal{M}_{\mathrm{in}}(b) \subset \mathcal{M}_{\mathrm{out}}$ guarantees the existence of a Type~I factor $\mathcal{N}$ interpolating between the two algebras, and hence a local tensor product decomposition of the Hilbert space. This is relevant to requirement~(2): in the presence of a split, the ``annular region'' between the inner and outer algebras has enough room to localize unitary implementers.

However, the split property does not by itself produce the transport of requirement~(1). It tells you that the relative position of $\mathcal{M}_{\mathrm{in}}(b)$ inside $\mathcal{M}_{\mathrm{out}}$ is ``standard'' in a measure-theoretic sense, but it does not provide a canonical family of automorphisms mapping $\mathcal{M}_*$ to $\mathcal{M}_{\mathrm{in}}(b)$ as $b$ varies. That family must come from an external source: either a symmetry group acting on the net, or an explicit transportability assumption.

Nor does the split property address requirement~(3). The $C^1$ regularity of the pulled-back state family and its tangent compatibility with the modular Hamiltonian variation are analytic conditions that depend on the smoothness of whatever generates the transport, not on the split structure.

In summary: split property assists with unitary implementability but does not generate transport or regularity. It is a facilitating condition, not a sufficient one.

\subsection{Diff-Covariant Nets}\label{diff-covariant-nets}

The cleanest sufficient condition for $(LImp^1)$ is diffeomorphism covariance of the net.

\begin{theorem}[Diff-covariance implies $(LImp^1)$]\label{thm:diff-limp}
Let $I \mapsto \mathcal{A}(I)$ be a local net on $S^1$ (or on the half-line) that is covariant under a strongly continuous positive-energy representation of $\mathrm{Diff}(S^1)$ (respectively, of orientation-preserving diffeomorphisms of the half-line fixing the boundary). Then $(LImp^1)$ holds for the cross-ratio family $\{\mathcal{M}_{\mathrm{in}}(b)\}_{b \in (0,1)}$ inside $\mathcal{M}_{\mathrm{out}}$.
\end{theorem}

\begin{proof}
Fix a reference point $b_0 \in (0,1)$ and set $\mathcal{M}_* = \mathcal{M}_{\mathrm{in}}(b_0)$. For each $b$ in a neighborhood of $b_0$, choose a smooth family of diffeomorphisms $f_b$ satisfying $f_{b_0} = \mathrm{id}$, $f_b(b_0) = b$, $f_b(1) = 1$, with support contained in the outer interval.

\emph{Transport.} The diffeomorphism covariance of the net gives $f_b(\mathcal{M}_*) = \mathcal{M}_{\mathrm{in}}(b)$, and since $\mathrm{supp}(f_b)$ is contained in the outer interval, the corresponding automorphism $\alpha_b$ restricts to an automorphism of $\mathcal{M}_{\mathrm{out}}$. This establishes requirement~(1).

\emph{Unitary implementability.} The positive-energy representation provides unitaries $U(f_b)$ implementing $\alpha_b$. Since $f_b$ is supported in the outer interval, $U(f_b) \in \mathcal{M}_{\mathrm{out}}$ by Haag duality (or by the localization properties of the Diff representation). Set $W_b := U(f_b)$. This establishes requirement~(2).

\emph{$C^1$ regularity.} The map $b \mapsto f_b$ is smooth by construction. The strong continuity of the Diff representation ensures that $b \mapsto U(f_b)$ is strongly continuous, and the smoothness of $f_b$ upgrades this to $C^1$ regularity of the pulled-back state family $b \mapsto \omega(W_b A W_b^*)$.

\emph{Tangent compatibility.} The infinitesimal generator of the family $U(f_b)$ at $b = b_0$ is a smeared stress-energy tensor supported in the outer interval. By the relation between stress-energy deformations and modular Hamiltonian variations in conformal field theory~\cite{CasiniHuertaMyers2011}, this generator represents the modular Hamiltonian variation $\dot{K}_{b_0}$ via the Araki correspondence. This establishes requirement~(3).

\emph{Nontriviality.} The diffeomorphisms $f_b$ are not Möbius transformations (they fix boundary points that a Möbius transformation would move). Therefore the pulled-back state family is nonconstant, and the QFIM is nonzero.
\end{proof}

This theorem upgrades the Virasoro worked example of Section~\ref{worked-example-ii-the-virasoro-net} from a model verification to an instance of a general result: any Diff-covariant conformal net automatically satisfies $(LImp^1)$.

\subsection{Transportability-Based Sufficiency}\label{transportability-sufficiency}

For nets that are not Diff-covariant, a second sufficient condition can be formulated by combining an explicit transportability hypothesis with the injectivity of the outer algebra.

\begin{definition}[Local transportability]\label{def:local-transport}
The cross-ratio family $\{\mathcal{M}_{\mathrm{in}}(b)\}$ is \emph{locally transportable} inside $\mathcal{M}_{\mathrm{out}}$ if, for each $b_0 \in (0,1)$, there exists a neighborhood $U_{b_0}$ and a family of automorphisms $\alpha_b \in \mathrm{Aut}(\mathcal{M}_{\mathrm{out}})$, $b \in U_{b_0}$, with $\alpha_{b_0} = \mathrm{id}$ and $\alpha_b(\mathcal{M}_*) = \mathcal{M}_{\mathrm{in}}(b)$.
\end{definition}

\begin{theorem}[Transportability-based sufficiency for $(LImp^1)$]\label{thm:transport-limp}
Suppose:
\begin{enumerate}
\item the cross-ratio family is locally transportable inside $\mathcal{M}_{\mathrm{out}}$;
\item $\mathcal{M}_{\mathrm{out}}$ is an injective (equivalently, hyperfinite) Type~III$_1$ factor;
\item the transport family $b \mapsto \alpha_b$ is $C^1$ in the sense that $b \mapsto \omega(\alpha_b(A))$ is $C^1$ for each $A \in \mathcal{M}_*$, with tangent compatible with the modular Hamiltonian variation.
\end{enumerate}
Then $(LImp^1)$ holds.
\end{theorem}

\begin{proof}
Requirement~(1) of $(LImp^1)$ is exactly hypothesis~(1). For requirement~(2): by hypothesis~(2) and the Connes theorem~\cite{Takesaki2003} ($\mathrm{Out}(\mathcal{M}) = \{e\}$ for injective Type~III$_1$ factors), every automorphism of $\mathcal{M}_{\mathrm{out}}$ is inner. Hence $\alpha_b = \mathrm{Ad}(W_b)$ for some $W_b \in \mathcal{U}(\mathcal{M}_{\mathrm{out}})$. Requirement~(3) is hypothesis~(3).
\end{proof}

\begin{remark}
If the inclusion $\mathcal{M}_{\mathrm{in}}(b) \subset \mathcal{M}_{\mathrm{out}}$ additionally satisfies the split property, then the Type~I interpolation allows $W_b$ to be chosen with support localized in the annular region between the two algebras. This localization can improve the analytic control needed for the $C^1$ regularity of hypothesis~(3), but it is not required for the theorem as stated.
\end{remark}

It is important to be clear about what this theorem does and does not achieve. Hypotheses~(1) and~(3) — transportability and $C^1$ tangent compatibility — are already the two main substantive components of $(LImp^1)$ itself. The theorem's contribution is to identify the precise algebraic condition (injectivity of the outer factor) that bridges the gap between ``transport exists as an automorphism'' and ``transport is unitarily implemented inside the outer algebra.'' For nets whose local algebras are known to be injective Type~III$_1$ factors — which includes all local algebras in standard AQFT models — this bridge is automatic.

\subsection{Status of the Three Phase 2 Conditions}\label{status-phase2-conditions}

The results of this section change the status of $(LImp^1)$ as follows:

\begin{itemize}
\item For Diff-covariant nets, $(LImp^1)$ is fully discharged by Theorem~\ref{thm:diff-limp}. No additional assumptions are needed.
\item For non-Diff-covariant nets with injective Type~III$_1$ local algebras, $(LImp^1)$ reduces to local transportability plus $C^1$ regularity (Theorem~\ref{thm:transport-limp}). The remaining burden is to verify that the cross-ratio family admits a smooth transport, which is a geometric property of the net rather than an abstract axiom.
\item $(CR\text{-}nd)$ and $(SFa)$ remain model-level verifications. They depend on dynamical and spectral properties (positivity of the stress-tensor two-point function, injectivity of the scale map) that are not captured by structural regularity conditions.
\end{itemize}

The net effect is that, for Diff-covariant nets, the Phase~2 conditionality is reduced from three independent axioms to two model-level verifications. For non-Diff-covariant nets, the situation is more modest: $(LImp^1)$ is decomposed and partially discharged, but the transportability and regularity inputs remain substantive.

\section{Verification: Two Worked
Examples}\label{verification-two-worked-examples}

The abstract Phase 2 result is deliberately conditional. Of the three
conditions, $(LImp^1)$ now has abstract sufficient conditions
(Section~\ref{sufficient-conditions-limp}): it is fully discharged for
Diff-covariant nets, and for other nets the remaining burden is isolated
to transportability and regularity. The remaining conditions $(CR\text{-}nd)$ and $(SFa)$
are verified in concrete models. This section provides that
verification in two structurally distinct examples.

The examples are chosen to do different kinds of work. The
Rindler/free-field chiral half-line model shows that the Phase 2 package
is realized in the same physical family that motivated much of the
original HAFF intuition. The Virasoro-net example then shows that the
construction is not a free-field accident and, in particular, that the
chart axiom becomes especially natural once Diff\((S^1)\) covariance is
built in.

\subsection{Worked Example I: Rindler / Free-Field Chiral
Half-Line}\label{worked-example-i-rindler-free-field-chiral-half-line}

The first example must be formulated carefully. The correct Phase 2
object is not the family of translated Rindler wedges. That family is
useful as a sanity check for the Phase 1 half-sided modular structure,
but it is the wrong object for the cross-ratio construction because
global translations leave the vacuum invariant and thus produce a
constant pulled-back state family. The correct model is the chiral
half-line realization associated with the Rindler horizon.

We therefore take: \[
\mathcal{M}_{\mathrm{out}}=\mathcal{A}(0,\infty),
\qquad
\mathcal{M}_{\mathrm{in}}(b)=\mathcal{A}(b,1),
\qquad
b\in(0,1),
\] in the vacuum representation of a chiral free-field net. The
canonical modular scale datum is then \[
K_b:=-\ln\Delta_{\mathcal{M}_{\mathrm{in}}(b),\Omega}.
\]

\subsubsection{\texorpdfstring{Verification of
\((LImp^1)\)}{Verification of (LImp\^{}1)}}\label{verification-of-limp1}

In this model, the required local chart is supplied by the local
diffeomorphism covariance of the free-field chiral net. Specifically, we
assume covariance under orientation-preserving diffeomorphisms of the
half-line that fix the boundary, implemented by unitaries in the
corresponding local algebras. Fix a reference point
\(b_0\in(0,1)\) and set \[
\mathcal{M}_*:=\mathcal{A}(b_0,1).
\] Choose a smooth local family of diffeomorphisms \[
f_b:(0,\infty)\to(0,\infty),
\qquad
b\in U_{b_0},
\] such that

\begin{enumerate}
\def\labelenumi{\arabic{enumi}.}
\tightlist
\item
  \(f_{b_0}=\mathrm{id}\);
\item
  \(f_b(1)=1\);
\item
  \(f_b(b_0)=b\);
\item
  the support of \(f_b\) is contained in the outer half-line.
\end{enumerate}

This covariance gives unitaries \[
W_b:=U(f_b)\in\mathcal{U}(\mathcal{M}_{\mathrm{out}})
\] with \[
W_b\mathcal{M}_*W_b^*=\mathcal{M}_{\mathrm{in}}(b).
\] Because the deformation is genuinely boundary-fixed rather than
Möbius, the family is nontrivial; because \(b\mapsto f_b\) is smooth,
the pulled-back state family is \(C^1\); and because the infinitesimal
generator is represented by the stress tensor, the tangent of the state
family is compatible with the modular Hamiltonian variation. Thus
\((LImp^1)\) holds in the model.

\subsubsection{\texorpdfstring{Verification of
\((CR\text{-}nd)\)}{Verification of (CR\textbackslash text\{-\}nd)}}\label{verification-of-crtext-nd}

In a vacuum chiral conformal theory, the modular Hamiltonian of an
interval \((a,b)\) has the local CHM form~\cite{CasiniHuertaMyers2011} \[
K_{(a,b)}
=
2\pi\int_a^b
\frac{(x-a)(b-x)}{b-a}\,
T_{00}(x)\,dx.
\] Specializing to the present family gives \[
K_b
=
2\pi\int_b^1
\frac{(x-b)(1-x)}{1-b}\,
T_{00}(x)\,dx,
\] and differentiation yields \[
\partial_b K_b
=
-2\pi\int_b^1
\frac{(1-x)^2}{(1-b)^2}\,
T_{00}(x)\,dx.
\] The corresponding Araki-Hessian quadratic form is positive and
nonzero, so the scale shadow is non-degenerate. More explicitly, one
obtains \[
h_b
=
C_R(1-b)^{-2}db^2,
\qquad
C_R>0.
\] Hence \((CR\text{-}nd)\) is satisfied.

\subsubsection{\texorpdfstring{Verification of
\((SFa)\)}{Verification of (SFa)}}\label{verification-of-sfa}

In this model, faithfulness reduces to a one-dimensional injectivity
statement. Since \[
h_b=C_R(1-b)^{-2}db^2
\] is strictly monotone in \(b\), the metric profile determines the
parameter \(b\), and hence determines the corresponding modular datum
\(K_b\) up to the trivial equivalences already allowed in the definition
of \((SFa)\). Therefore the scale map is faithful in the required sense.

The upshot is that the entire Phase 2 package is realized in the
Rindler/free-field chiral half-line model: \[
(\mathrm{LImp}^1),\quad
(\mathrm{CR\text{-}nd}),\quad
(\mathrm{SFa}).
\]

\subsection{Worked Example II: The Virasoro
Net}\label{worked-example-ii-the-virasoro-net}

The second example shows that the preceding mechanism is not tied to
free fields. Let \[
I\longmapsto \mathcal{A}_{\mathrm{Vir},c}(I)
\] be the Virasoro net at central charge \(c\), with vacuum vector
\(\Omega\). As in the first example, we work in the standard half-line
gauge: \[
\mathcal{M}_{\mathrm{out}}=\mathcal{A}_{\mathrm{Vir},c}(0,\infty),
\qquad
\mathcal{M}_{\mathrm{in}}(b)=\mathcal{A}_{\mathrm{Vir},c}(b,1).
\]

\subsubsection{\texorpdfstring{Verification of
\((LImp^1)\)}{Verification of (LImp\^{}1)}}\label{verification-of-limp1-1}

What was model-dependent in the first example becomes built in here. The
Virasoro net is Diff\((S^1)\)-covariant, so \((LImp^1)\) follows
directly from Theorem~\ref{thm:diff-limp}. Concretely, one can choose a
smooth boundary-fixed family \[
f_b\in\mathrm{Diff}(S^1),
\qquad
b\in U_{b_0},
\] with \[
f_{b_0}=\mathrm{id},
\qquad
f_b(b_0)=b,
\qquad
f_b(1)=1,
\] and support contained in the outer half-line. The positive-energy
Diff\((S^1)\) representation yields unitaries \[
W_b:=U(f_b)\in\mathcal{U}(\mathcal{M}_{\mathrm{out}})
\] such that \[
W_b\mathcal{A}(b_0,1)W_b^*=\mathcal{A}(b,1).
\] The nontriviality, smoothness, and tangent compatibility required by
\((LImp^1)\) follow immediately from the non-Möbius character of
\(f_b\), the smoothness of the Diff\((S^1)\) action, and the fact that
the infinitesimal action is generated by the Virasoro stress tensor.
Thus \((LImp^1)\) holds here in a cleaner and more canonical way than in
the free-field model.

\subsubsection{\texorpdfstring{Verification of
\((CR\text{-}nd)\)}{Verification of (CR\textbackslash text\{-\}nd)}}\label{verification-of-crtext-nd-1}

The interval modular Hamiltonian again has the CHM form \[
K_b
=
2\pi\int_b^1
\frac{(x-b)(1-x)}{1-b}\,
T(x)\,dx,
\] where \(T\) is now the Virasoro stress tensor. Differentiating gives
the same structural formula as before: \[
\partial_b K_b
=
-2\pi\int_b^1
\frac{(1-x)^2}{(1-b)^2}\,
T(x)\,dx.
\] Using the Virasoro two-point function \[
\langle T(x)T(y)\rangle\sim \frac{c/2}{(x-y)^4},
\] one obtains, after standard regularization, a positive finite
coefficient \[
\mathbf{Q}_{\omega,b}(\partial_b K_b,\partial_b K_b)
=
\frac{C_{\mathrm{Vir}}(c)}{(1-b)^2},
\qquad
C_{\mathrm{Vir}}(c)>0.
\] Hence \[
h_b
=
C_{\mathrm{Vir}}(c)(1-b)^{-2}db^2,
\] and \((CR\text{-}nd)\) holds.

\subsubsection{\texorpdfstring{Verification of
\((SFa)\)}{Verification of (SFa)}}\label{verification-of-sfa-1}

Exactly as in the first example, the scale profile is strictly injective
in the cross-ratio parameter. Therefore the full scale family determines
\(b\), and hence determines the local modular datum up to the same
trivial equivalences. Thus \((SFa)\) is satisfied in the Virasoro model
as well.

\subsection{Comparison and Universal
Scaling}\label{comparison-and-universal-scaling}

The significance of the two examples lies less in their individual
computations than in what they share.

\begin{enumerate}
\def\labelenumi{\arabic{enumi}.}
\tightlist
\item
  Both models realize the same abstract Phase 2 package: \[
  (\mathrm{LImp}^1)+(\mathrm{CR\text{-}nd})+(\mathrm{SFa}).
  \]
\item
  Both produce the same scale profile: \[
  h_b=C(1-b)^{-2}db^2.
  \]
\item
  The model dependence enters only through the coefficient:

  \begin{itemize}
  \tightlist
  \item
    in the free-field model through the normalization of the
    stress-tensor sector;
  \item
    in the Virasoro model through the central charge \(c\).
  \end{itemize}
\end{enumerate}

This is the strongest available evidence, within the present paper, that
the scale shadow is not a model-specific artifact. What appears to be
universal is the cross-ratio scaling law itself, which is fixed by the
CHM structure of interval modular Hamiltonians. What remains
model-sensitive is only the overall constant.

From the point of view of the paper's philosophical thesis, this matters
because it shows that the full inseparability mechanism is not tied to
one special solvable example. The rigidity claim is abstract; the
examples show that the additional Phase 2 inputs are concretely
realizable in more than one structural setting, and that the theorem's
conditions arise naturally where one would most expect them to arise.

\section{Implications for the Hard
Problem}\label{implications-for-the-hard-problem}

For the purposes of this section, we take the mathematical results of
the previous sections as given. At the conformal level, Theorem 4.8
already shows that conformal geometry and conformal modular data are not
independently deformable. The later Phase 2 results strengthen this to
the full scale-sensitive level by tying the metric scale to the
canonical modular scale datum. The philosophical question is therefore no longer
whether one can attach a speculative interpretation to a formal analogy.
It is whether the standard formulation of the hard problem survives once
its key presupposition is written in mathematically precise terms.

\subsection{The Failure of the
Presupposition}\label{the-failure-of-the-presupposition}

Section 1 identified the hidden presupposition of the hard problem:
independent describability. Section 2 translated it into independent
deformability. The inseparability theorems of Sections 4 and 5 now
supply the answer.

At the Phase 1 level, Theorem 4.8 shows that conformal geometry and
conformal modular data are already locked together. At the full level,
Theorem 5.7 shows that, under \((LImp^1)\), \((CR\text{-}nd)\), and
\((SFa)\), the full geometric datum and the extended modular
data cannot be independently varied. Conditional on the identifications
defended in Section 2, the presupposition is not available at either
level.

The strongest claim justified by the theorem is therefore not that the
hard problem has been ``solved'' in the ordinary sense. It is that its
canonical production form is ill-posed. The question ``how does matter
produce mind?'' asks for a derivation across a divide that, in the
present framework, is not mathematically available as an independent
divide. One cannot first specify the geometric-physical side and only
afterwards ask how the first-person side appears. They are two
structurally inseparable readings of the same algebraic datum
\((\mathcal{A},\omega)\).

\subsection{Relocating the Explanatory
Gap}\label{relocating-the-explanatory-gap}

Once the production picture is rejected, a residual question remains.
Even if geometry and modular flow are inseparable, why should a given
modular invariant correspond to a given phenomenal quality? This is
legitimate, but it is not the original hard problem. It concerns the
interpretive reach of an observer embedded in the very structure being
described.

The proper location of the remaining gap is not between physical and
experiential structure as two externally juxtaposed domains. It lies
within the self-description limits of a finite agent modeled by
\((\mathcal{A},\omega)\). The incompleteness result of Paper G~\cite{Liu2026PaperG}
(which proves that no accessibility-based framework can fully
internalize its own accessibility conditions without either regress or
collapse, because the self-referential fixed point
\(\mathcal{A}^*\in\mathrm{Acc}(\mathcal{A}^*)\) generically does not
exist) shows
that a finite self-model cannot be fully closed under its own
self-referential demands. The agent cannot, from within the same
structure, provide a fully transparent answer to why one particular
modular pattern is lived as this particular qualitative articulation.

The explanatory gap thereby becomes structural rather than
technical~\cite{Levine1983}. It
is not a hole waiting for a better neuroscientific mechanism. It is a
horizon effect generated by self-reference and embedding, closer to a
Gödelian limitation than to pre-Newtonian ignorance of gravity.

This changes what counts as philosophical success. The goal is no longer
to show how an independently specified physical world outputs
experience. It is to understand why an embedded agent necessarily
encounters a limit in its attempt to completely internalize the relation
between structural invariants and lived articulation. The residual gap
is relocated and reclassified, not removed.

\subsection{Comparison with Existing
Positions}\label{comparison-with-existing-positions}

The closest philosophical relatives of the present framework are
positions that deny a brute separation between mind and world. Even so,
the relation is exact only up to a point. The following comparison is
necessarily compressed; each position listed has significant internal
diversity, and the characterizations below refer to the canonical or
most common formulation.

\begin{center}
\begin{tabularx}{0.98\linewidth}{@{}P{0.23\linewidth}Y@{}}
\toprule
Position & Relation to the present paper \\
\midrule
Physicalism / functionalism & Shares the desire for structural rigor,
but typically preserves an explanatory asymmetry: physical facts are
taken as independently specifiable, and consciousness must be shown to
follow. The present theorem denies that asymmetry in its standard
form. \\
Panpsychism & Avoids emergence by treating experience as basic, but
inherits the combination problem. The present framework does not posit
basic micro-experiences; it treats modular structure and geometry as
co-emergent readings of one algebraic organization. \\
Dual-aspect monism & The nearest metaphysical relative. The decisive
difference is that here the dual-aspect claim is not left at the level
of metaphysical slogan; it is tied to a concrete inseparability
theorem. \\
IIT & Also seeks a mathematically articulated bridge, but introduces an
additional quantity \(\Phi\) as a privileged measure. The present route
instead treats inseparability as internal to modular and geometric
structure, not as the effect of an extra consciousness functional. \\
Illusionism & Rejects the hard problem by denying that phenomenal
consciousness needs the explanatory role usually assigned to it. The
present framework does not deny experience. It denies the independent
specification presupposition that makes experience look like a
separately generated add-on. \\
Russellian monism & Very close in spirit insofar as it challenges the
exhaustiveness of third-person structure. But Russellian monism often
speaks of an intrinsic categorical basis behind structural physics. The
present framework is more austere: both geometry and first-person time
are readings of one algebraic structure, without appeal to a hidden
intrinsic substrate beyond that structure. \\
\bottomrule
\end{tabularx}
\end{center}

The upshot is that the present work should not be advertised as a simple
algebraic version of any single existing position. It is closest to a
theorem-backed dual-aspect view, but even that label is only
approximate. The real novelty lies in replacing a metaphysical balance
claim with a rigidity statement.

A structural correspondence with the Yogacara fourfold analysis is
recorded in Appendix~\ref{app:yogacara}.

\section{Broader Significance}\label{broader-significance}

The philosophical argument of the preceding sections is the primary
contribution of this paper. The mathematical and physical content,
however, has independent interest that is worth recording briefly.

\subsection{Mathematical Significance}\label{mathematical-significance}

The Wiesbrock-GLW programme~\cite{Wiesbrock1993,GLW1998} established that half-sided modular
inclusions reconstruct conformal symmetry. What remained open was the
passage from conformal skeleton to metric scale using only modular data.
The present paper isolates a precise conditional route to a full
geometric datum in 1+1 dimensions: the cross-ratio modular scale datum
\(\{K_b\}\) serves as the canonical non-gauge geometric modulus, and the
Araki relative entropy Hessian~\cite{Araki1976} serves as the Type III-compatible
QFIM that converts it into a Riemannian scale shadow. The route is
conditional on \((LImp^1)\), \((CR\text{-}nd)\), and \((SFa)\), all
three of which are verified in two structurally distinct models.
For Diff-covariant nets, Section~\ref{sufficient-conditions-limp}
further shows that \((LImp^1)\) is automatic, leaving only
\((CR\text{-}nd)\) and \((SFa)\) as model-level verifications.
In the
worked examples, where the cross-ratio parameter naturally parametrizes
the spatial slice, this datum assembles into the Lorentzian metric.

The universal scaling law \(h_b \sim (1-b)^{-2}db^2\), shared by both
worked examples, is a structural consequence of the Casini-Huerta-Myers
form~\cite{CasiniHuertaMyers2011} of interval modular Hamiltonians. It is not a free-field accident.
The model-dependent information enters only through the overall
constant, which encodes data such as the central charge.

\subsection{Physical Significance}\label{physical-significance}

The slogan ``it from qubit'' and the ER=EPR programme assert that
entanglement structure determines spacetime geometry~\cite{Maldacena1999,RyuTakayanagi2006,VanRaamsdonk2010,Maldacena2013}. These claims have
so far been tied to the AdS/CFT correspondence. The present result
establishes a version of the same thesis without assuming holographic
duality. The axiom package
selected by \((U3)\), together with the axiom package
\((U1)+(C)+(Amb)+(CN_\exists)+(LImp^1)+(CR\text{-}nd)+(SFa)\), makes
no reference to anti-de Sitter space, conformal field theory on a
boundary, or bulk-boundary correspondence. Inseparability is a property
of Type III\(_1\) algebraic structure under appropriate
regularity~\cite{Witten2018}, not
a special consequence of holography.

This also resolves two open issues from the earlier HAFF programme
(Paper D)~\cite{Liu2026PaperD}. The non-uniqueness of the map from algebraic deformations to
state deformations is absorbed by the cross-ratio datum and the chart
axiom, which together fix the relevant deformation direction. The
signature mismatch between the positive-definite QFIM and the Lorentzian
metric is resolved by restricting the QFIM to spatial slices and letting
the conformal skeleton supply the causal structure.

\section{Explicit Non-Claims and Open
Problems}\label{explicit-non-claims-and-open-problems}

\subsection{What This Paper Does Not
Claim}\label{what-this-paper-does-not-claim}

\begin{enumerate}
\def\labelenumi{\arabic{enumi}.}
\item
  It does not claim to have solved the hard problem of consciousness. It
  claims to have shown that the standard production formulation rests on
  a presupposition that is not available in the relevant mathematical
  framework.
\item
  It does not claim that modular flow ``is'' consciousness. The
  identification in Section 2 is a structural correspondence: modular
  flow plays the mathematical role that first-person temporal structure
  plays in the philosophical debate.
\item
  It does not claim that the inseparability theorem holds
  unconditionally for all Type III\(_1\) factors. Phase 1 is
  unconditional within the stated axiom class. Phase 2 is conditional on
  \((LImp^1)\), \((CR\text{-}nd)\), and \((SFa)\). For Diff-covariant nets,
  \((LImp^1)\) is discharged by Theorem~\ref{thm:diff-limp}; the remaining
  two conditions are verified in two models but not proved in full generality.
\item
  It does not claim that the Buddhist correspondence in Section 7.4
  participates in any proof. It is included as a structural comparison,
  not as evidence.
\item
  It does not claim that the results extend to dimensions higher than
  1+1. The GLW construction produces a Möbius net on \(S^1\).
  Higher-dimensional generalizations require additional HSMI directions
  and remain open.
\item
  It does not claim that consciousness is a quantum effect. The argument
  concerns the structural relationship between geometry and modular flow
  at the level of algebraic quantum theory, not a physical mechanism
  involving quantum coherence in biological systems.
\end{enumerate}

\subsection{Open Problems}\label{open-problems}

\begin{enumerate}
\def\labelenumi{\arabic{enumi}.}
\item
  Abstract sufficient conditions for \((CR\text{-}nd)\)
  and \((SFa)\). Section~\ref{sufficient-conditions-limp} provides sufficient conditions for \((LImp^1)\) (Diff-covariance or split-assisted transportability), but the remaining two conditions are still verified only in specific models. A general criterion for \((CR\text{-}nd)\), perhaps in terms of spectral properties of the stress-energy tensor, would further reduce the conditionality.
\item
  Foliation compatibility beyond 1+1 dimensions. In the present 1+1
  dimensional setting, the cross-ratio family is a single-parameter
  family and the assembly is straightforward. In higher dimensions, the
  scale data from different cross-ratio slices must satisfy nontrivial
  mutual consistency conditions. Formulating and verifying these
  conditions is an open problem tied to the higher-dimensional
  extension.
\item
  Higher-dimensional extension. In \(d > 1+1\), the conformal group is
  larger and the cross-ratio construction must be replaced by a
  higher-dimensional analogue. Whether the inseparability mechanism
  survives this extension is the most important structural open
  question.
\item
  Quantitative Paper G. The incompleteness result invoked in Section 7.2
  is currently qualitative. A quantitative version, bounding the
  self-description capacity of a finite observer in terms of modular
  invariants, would sharpen the relocation of the explanatory gap.
\item
  The relationship between the present framework and Integrated
  Information Theory. Both seek mathematically articulated accounts of
  consciousness. Whether the inseparability theorem implies constraints
  on \(\Phi\) or vice versa is unexplored.
\end{enumerate}

\appendix

\section{Yogacara Fourfold Correspondence}\label{app:yogacara}

The Yogacara fourfold analysis provides a structural correspondence for
the architecture that emerges in this paper. The correspondence does not
enter any proof; it is included as a disciplined interpretive comparison
that becomes newly intelligible once the mathematical theorem is in
place. Each position listed below has significant internal diversity
within the Buddhist philosophical tradition; the characterizations refer
to the canonical formulation.

{\def\LTcaptype{none}
\hbadness=10000
\begin{longtable}{P{0.14\linewidth}P{0.24\linewidth}P{0.54\linewidth}}
\toprule\noalign{}
Layer & Phenomenological role & Mathematical counterpart \\
\midrule\noalign{}
\endhead
\bottomrule\noalign{}
\endlastfoot
相分 & the presented object-side & emergent geometry
\(\mathrm{Geom}\) \\
见分 & the act-side of cognition & modular flow \(\sigma_t^\omega\) \\
自证分 & reflexive awareness & the self-calibrating informational loop
of T-DOME Paper~III~\cite{Liu2026TDOME_III} \\
证自证分 & the boundary of reflexivity & the incompleteness limit of
Paper~G~\cite{Liu2026PaperG} \\
\end{longtable}
}

The structural fit is that object-side and act-side are not heterogeneous
substances but internally linked aspects of a single cognitive event;
reflexivity does not continue indefinitely into an infinite hierarchy,
but closes through a higher-order boundary condition. This is exactly
the pattern suggested by the inseparability framework: geometry and
modular flow are inseparable at the base level, and reflexive closure is
marked by incompleteness rather than by defect.
 
\newpage
%
\addcontentsline{toc}{section}{References}

\begin{thebibliography}{99}

\bibitem{Araki1976}
H.~Araki,
\emph{Relative entropy of states of von Neumann algebras},
Publ.\ Res.\ Inst.\ Math.\ Sci.\ Kyoto Univ.\ \textbf{11}, 809--833 (1976).

\bibitem{Borchers1992}
H.-J.~Borchers,
\emph{The CPT theorem in two-dimensional theories of local observables},
Commun.\ Math.\ Phys.\ \textbf{143}, 315--332 (1992).

\bibitem{CasiniHuertaMyers2011}
H.~Casini, M.~Huerta, and R.~C.~Myers,
\emph{Towards a derivation of holographic entanglement entropy},
JHEP \textbf{1105}, 036 (2011),
arXiv:1102.0440.

\bibitem{Chalmers1996}
D.~J.~Chalmers,
\emph{The Conscious Mind: In Search of a Fundamental Theory},
Oxford University Press (1996).

\bibitem{GLW1998}
R.~Guido, R.~Longo, and H.-W.~Wiesbrock,
\emph{Extensions of conformal nets and superselection structures},
Commun.\ Math.\ Phys.\ \textbf{192}, 217--244 (1998).

\bibitem{Haag1996}
R.~Haag,
\emph{Local Quantum Physics: Fields, Particles, Algebras},
2nd ed., Springer (1996).

\bibitem{Levine1983}
J.~Levine,
\emph{Materialism and qualia: The explanatory gap},
Pacific Philosophical Quarterly \textbf{64}, 354 (1983).

\bibitem{Liu2026PaperD}
S.~Liu,
\emph{Gravitational Phenomena as Emergent Properties of Observable Algebra Selection: A Structural Analysis},
Zenodo (2026), DOI: 10.5281/zenodo.18388881.

\bibitem{Liu2026PaperG}
S.~Liu,
\emph{Structural Limits of Unification: Accessibility, Incompleteness, and the Necessity of a Final Cut},
Zenodo (2026), DOI: 10.5281/zenodo.18402907.

\bibitem{Liu2026TDOME_III}
S.~Liu,
\emph{Fisher Information Geometry and the Thermodynamic Cost of Self-Referential Calibration},
Zenodo (2026), DOI: 10.5281/zenodo.18591771.

\bibitem{Maldacena1999}
J.~M.~Maldacena,
\emph{The large N limit of superconformal field theories and supergravity},
Int.\ J.\ Theor.\ Phys.\ \textbf{38}, 1113 (1999).

\bibitem{Maldacena2013}
J.~Maldacena and L.~Susskind,
\emph{Cool horizons for entangled black holes},
Fortsch.\ Phys.\ \textbf{61}, 781 (2013),
arXiv:1306.0533.

\bibitem{RyuTakayanagi2006}
S.~Ryu and T.~Takayanagi,
\emph{Holographic derivation of entanglement entropy from AdS/CFT},
Phys.\ Rev.\ Lett.\ \textbf{96}, 181602 (2006).

\bibitem{Takesaki2003}
M.~Takesaki,
\emph{Theory of Operator Algebras II},
Encyclopaedia of Mathematical Sciences \textbf{125}, Springer (2003).

\bibitem{Unruh1976}
W.~G.~Unruh,
\emph{Notes on black-hole evaporation},
Phys.\ Rev.\ D \textbf{14}, 870 (1976).

\bibitem{VanRaamsdonk2010}
M.~Van~Raamsdonk,
\emph{Building up spacetime with quantum entanglement},
Gen.\ Rel.\ Grav.\ \textbf{42}, 2323 (2010).

\bibitem{Wiesbrock1993}
H.-W.~Wiesbrock,
\emph{Half-sided modular inclusions of von Neumann algebras},
Commun.\ Math.\ Phys.\ \textbf{157}, 83 (1993).

\bibitem{Witten2018}
E.~Witten,
\emph{APS Medal for Exceptional Achievement in Research: Invited article on entanglement properties of quantum field theory},
Rev.\ Mod.\ Phys.\ \textbf{90}, 045003 (2018),
arXiv:1803.04993.

\end{thebibliography}
 
\end{document}
